Nach der Regel von Dulong-Petit ist die molare Wärme fester Elemente bei hohen Temperaturen $C = 3R$ (FoTaBe S.~171). Mit der molaren Masse von Palladium $M = \MPdO$ (FoTaBe S.~229) ist die spezifische Wärme
\begin{align}
 c &= \frac{C}{M} = \cPdF = \frac{3\cdot\ncR}{\MPd} = \cPd \approx \result{\cPdP{0}{3}}
\end{align}
